\scanpage{18}
\mspage{(16)}

\unsure{complexes} de De Rham à p.d. \inserted{(ou même $D^{\bullet\bullet}(\mathcal X_*,\mathbb Z)$)} \unsure{sufft|suffit} déjà pour récupérer toute la structure mult. à p.d. coh. de $\mathcal X_*$ \inserted{(sous réserve des conditions de finitude bien sûr)}. Dans le cas contraire, ce serait cette dernière qui serait le candidat \inserted{alg. ``linéaire''} naturel pour exprimer le type d'homotopie $\mathcal X_*$ (du moins si $\mathcal X_*$ \inserted{connexe et} simplement connexe, et en passant à une cat. dérivée convenable bien sûr).

\mspage{17}
